% Immutable cited-input supplement. Target proofs remain embargoed.
% Target: night:P4.22; tag 02WU.
% Frozen source SHA256 74497bfe327fae25c9cdcb17334b3bf42a9793a9df7c0a4404d97568f2f760f0.
% CompendiumRef is resolved against the cited S6_09 source edition, not silently identified with the current series statement.

% Exact cited input: series:II / night:P4.13
\begin{proposition}[P4.13]\hypertarget{tag-02WK}{}\label{night:P4.13}\label{stmt:P4.13}


Let a proper holomorphic degeneration of two-dimensional complex
tori have smooth total threefold and central fibre $mS$ with
irreducible reduced support. If $A$ is the monodromy on $H_1$ of
a nearby torus, then
\[
 m e(S_{\mathrm{reg}})=\det(I-A^m),\qquad
 0\le e(S_{\mathrm{reg}})\le16/m.
\]
This assertion requires no Betti-number hypothesis on a compactification.
\end{proposition}

% Exact cited input: series:II / night:P4.15
\begin{proposition}[P4.15]\hypertarget{tag-02WM}{}\label{night:P4.15}\label{stmt:P4.15}


Let $mS$ be an irreducible singular fibre in the setting of 4.2.
Suppose its normalization $T$ is smooth and $d\nu:T_T\to\nu^*T_X$
has rank two away from finitely many points. Let $\ell$ be the
length of its maximal-minor scheme. Then
\[
 e(T)=K_T^2-\ell.
\]
The alternatives for $T$ are a rational surface with
$e(T)=6-\ell/2$ and $0\leq\ell\leq6$, a minimal ruled surface
over an elliptic curve with $\ell=0$, or a Hopf surface with
$\ell=0$. In either of the last two alternatives $e(S)=0$.
The reduced conductor and singular locus are disjoint unions of
smooth elliptic curves, and their normalization maps are finite
and unramified.
\end{proposition}

% Exact cited input: series:II / night:P4R.1
\begin{proposition}[P4R.1]\hypertarget{tag-02Y2}{}\label{night:P4R.1}\label{stmt:P4R.1}


For a proper holomorphic degeneration of complex two-tori with
smooth total threefold and central fibre $mS$, where $S$ is
irreducible and reduced, the group $\pi_1(S)$ has an abelian
normal subgroup with cyclic quotient of order dividing $m$.
\end{proposition}

% Exact cited input: series:II / night:P4R.3
\begin{proposition}[P4R.3]\hypertarget{tag-02YF}{}\label{night:P4R.3}\label{stmt:P4R.3}


In the standing torus-fibration setting of this section, let an irreducible singular fibre
support $S$ have smooth normalization $\nu:T\to S$ whose
minimal model is ruled over a curve of genus $g\ge2$.
Write $D$ for the conductor divisor on $T$ and
$C=\operatorname{Sing}(S)_{\rm red}$.
The inclusion $D_{\rm red}\to T$ induces an isomorphism on
fundamental groups, and the normalization pushout induces
$\pi_1(S)\simeq\pi_1(C)$. Moreover,
\[
 e(C)\ge0,\qquad e(S)>0.
\]
The assertion includes conductor components on different
ruling fibres that map onto the same curve in $C$.
\end{proposition}

